\section{Uniform one-pass computation with charged bit-space}\label{sec:streaming64}
The preceding label-width invariant permits arbitrary real transition
tables. A lower bound for that invariant also constrains finite-bit
algorithms, but an upper bound for it need not be implementable at a
comparable cost. We now give a separate uniform result in which all
integer storage and arithmetic are charged. Throughout this section the
physical alphabet is exactly $\mathcal A_6$ in
Theorem~\ref{thm:noidle62}; in particular, no idle command is added.
All logarithms are natural except those explicitly denoted by $\log_2$.

\subsection{Computational convention and the space theorem}
A numerical streaming algorithm has a fixed finite program. It reads
parameters $N\ge2$ and $L\ge2$, followed by $N$ commands, in one pass.
A consumed command cannot be consulted again. No input-dependent
information may be stored in a readable output tape, an external
clock, or an uncharged address. After the last command it returns a
rational matrix $\widehat F(w)\in\mathcal D_2$. For randomized
algorithms every returned matrix is required to be legal, and our
correctness condition is the weaker, mean condition
\begin{equation}\label{eq:streammean64}
 \sup_{w\in\mathcal A_6^N}
 \bigl\|\E\widehat F(w)-D(A_we_1)\bigr\|_{\rm F}\le2^{-L}.
\end{equation}
Internal computation terminates almost surely between successive
commands and at the output stage. Randomness is fresh; any retained
random seed is part of the configuration.

Writable bit-space includes the data registers, temporary integer
storage, retained parameter information, work-tape addresses and
control information. Equivalently, one may charge a binary encoding of
the complete private configuration. To relate these conventions,
a fixed program with $s$ work bits and a fixed number of work heads
has at most $C(s+1)^d2^{cs}$ configurations, for constants depending
only on its machine convention. The prescribed position of the
one-pass input head at the $t$th command boundary contains no
input-dependent information. A private address that does contain such
information must be counted. The parameters are not an advice table;
they contain only the two integers $N,L$.

\begin{theorem}[Uniform streaming space]\label{thm:streamspace64}
There is one deterministic numerical streaming algorithm for the six
rational rotations such that, for every $N\ge2$ and $L\ge2$, every
output is legal, the pointwise error is at most $2^{-L}$, and the
writable bit-space is
\begin{equation}\label{eq:spaceupper64}
 O\bigl(\min\{N,L+\log_2(N+1)\}\bigr).
\end{equation}
It uses no real advice, no random sampling, and no oracle arithmetic.
Writing $b_* =\min\{N,L+\lceil\log_2N\rceil+2\}$, its command updates
can be performed in $O(b_*)$ bit operations each with $O(b_*)$ scratch
space. Parameter scanning is charged separately by its input length;
a binary rational output uses $O(b_*)$ bits.

Conversely, there is an absolute $c>0$ such that every one-pass
algorithm satisfying \eqref{eq:streammean64} and having at most $2^B$
complete private configurations at each command boundary satisfies
\begin{equation}\label{eq:spacelower64}
 B\ge c\min\{N,L+\log_2(N+1)\}.
\end{equation}
There is no running-time restriction in this converse. It applies even
to arbitrary horizon-dependent stochastic updates. Consequently every
fixed-program randomized bit algorithm has the same lower order of
writable space, with machine-dependent constants. Deterministic and
randomized uniform space therefore have the common order
$\Theta(\min\{N,L+\log_2(N+1)\})$ for this numerical task.
\end{theorem}
The spectral input for the converse is the full spherical Koopman norm
$\sqrt5/3$ of the six-letter average, not the norm of its action on
$\R^3$. The upper algorithm uses only rational orthogonality and does
not need a spectral assumption. The theorem is a consequence of the
entropy occupation law together with an explicit finite-arithmetic
construction, not a new proof of the LPS spectral theorem.

\subsection{From complete configurations to stochastic width}
\begin{lemma}[Boundary reduction]\label{lem:boundary64}
If a one-pass algorithm has at most $k$ private configurations at every
command boundary and satisfies mean error $\varepsilon$, its boundary
laws define a clocked numerical realization of peak width at most $k$
and the same error.
\end{lemma}
\begin{proof}
Take the configurations at boundary $t$ as $S_t$. Include every piece
of retained information, whether or not it is called a data register.
Starting from a boundary configuration and a supplied next command,
integrate over all internal computation and fresh randomness up to the
next boundary. Almost-sure termination gives a stochastic row of total
mass one. The row depends only on that configuration and command,
with the fixed parameters and boundary index permitted as public
parameters. Consumed input and old random bits are inaccessible unless
recorded in the configuration. This gives precisely the Markov rows
of Definition~\ref{def:machine54}.

At the terminal boundary, replace any randomized output computation
by its conditional mean matrix. This mean is in $\mathcal D_2$ by
convexity and compactness, and is independent of the discarded input
conditional on the complete configuration. Thus the numerical decoder
is legal and has exactly the mean output in
\eqref{eq:streammean64}. No running-time or exact-sampling assumption
is introduced in this reduction.
\end{proof}

\begin{proposition}[A finite bit--accuracy obstruction]\label{prop:bitaccuracy64}
For the six-command experiment, suppose a correct realization has at
most $2^B$ complete configurations at every positive cut and error
$0\le\varepsilon<1/\sqrt2$. Put
\[
 a_B=1+6\log(12\cdot2^B),\qquad c_B=648\pi^2\,2^B.
\]
Then
\begin{equation}\label{eq:finitebits64}
 N\le a_B+c_B\frac{\sqrt2\varepsilon}{1-\sqrt2\varepsilon}.
\end{equation}
In particular, if $d=N-a_B>0$, then
\begin{equation}\label{eq:invertedbits64}
 \varepsilon\ge\frac{d}{\sqrt2(c_B+d)}.
\end{equation}
At zero error the exact profile gives the stronger bound
$B\ge N\log_2 5$.
\end{proposition}
\begin{proof}
Apply \eqref{eq:noidleoccupation62} to all $N$ positive cuts with
$k=2^B$ and $\zeta=1-\sqrt2\varepsilon$. This gives
\eqref{eq:finitebits64}, including its limiting zero-error case.
Solving $d\le c_B\sqrt2\varepsilon/(1-\sqrt2\varepsilon)$ gives
\eqref{eq:invertedbits64}. For the last assertion use the exact
terminal width $5^N$ in Theorem~\ref{thm:noidle62} instead of the
weaker entropy estimate.
\end{proof}

\begin{proof}[Lower bound in Theorem~\ref{thm:streamspace64}]
Let $k$ be the maximum number of boundary configurations. Since
$\varepsilon=2^{-L}\le1/4$, Proposition~\ref{prop:bitaccuracy64}, or
\eqref{eq:noidleoccupation62} with $k$, implies, for absolute positive
constants $C_0,C_1,C_2$,
\begin{equation}\label{eq:bitdichotomy64}
 N\le C_0+C_1\log k+C_2\varepsilon k.
\end{equation}
For $N\ge2C_0$, either $\log k\ge N/(4C_1)$ or
$k\ge N/(4C_2\varepsilon)$. Taking logarithms and enlarging an
absolute additive constant yields
\[
 \log_2 k\ge c_0\min\{N,\log_2(N/\varepsilon)\}-C_3.
\]
The additive constant can be absorbed uniformly. Indeed $k\ge2$ for
all $N\ge2$ in the stated error range: the all-$A_1$ word ends at
$e_1$, whereas $N-1$ copies of $A_1$ followed by $A_2$ end at
$A_2e_1=(-3/5,0,-4/5)$. Their density matrices have Frobenius
distance $\sqrt{8/5}>1/2$. A single constant mean decoder cannot be
within $1/4$ of both. If
$\min\{N,\log_2(N/\varepsilon)\}$ is below a fixed large constant,
$\log_2k\ge1$ suffices; above that constant it absorbs $C_3$ in the
preceding estimate. The finitely bounded range $N<2C_0$ is handled
in the same way. Since $k\le2^B$, this proves
\eqref{eq:spacelower64}. Finally the fixed-program configuration count
stated above gives $\log_2k=O(s+1)$, so the same conclusion holds
for work-tape space. This argument neither assumes uniform transition
tables nor assumes that an individual sampled decoder approximates
the physical target.
\end{proof}

\subsection{A legality-preserving integer algorithm}
For $b\ge0$, define coordinatewise truncation toward zero by
\[
 T_b(x)_i=2^{-b}\operatorname{sgn}(x_i)
                          \lfloor2^b|x_i|\rfloor.
\]
Unlike rounding to the nearest grid point, this operation preserves
the closed Euclidean unit ball.
\begin{lemma}[Directed rounding]\label{lem:directedround64}
For every $x\in\R^d$,
\[
 \norm{T_b(x)}\le\norm x,\qquad
 \norm{T_b(x)-x}\le\sqrt d\,2^{-b}.
\]
If $R_1,\ldots,R_N$ are orthogonal, $x_t=R_tx_{t-1}$ with
$\norm{x_0}\le1$, and $v_0=x_0$ is on the $2^{-b}$ grid, then
$v_t=T_b(R_tv_{t-1})$ satisfies
\begin{equation}\label{eq:roundorbit64}
 \norm{v_t}\le1,\qquad \norm{v_t-x_t}\le t\sqrt d\,2^{-b}.
\end{equation}
\end{lemma}
\begin{proof}
Each absolute coordinate decreases under truncation and changes by
less than $2^{-b}$. Squaring and summing gives the first two estimates.
Orthogonality and the triangle inequality give
$\norm{v_t-x_t}\le\sqrt d\,2^{-b}+\norm{v_{t-1}-x_{t-1}}$.
Induction proves both assertions in \eqref{eq:roundorbit64}.
\end{proof}

Write $A_a=M_a/5$, where $M_a$ is one of the six fixed integer
matrices from \eqref{eq:lpsinput61} and their transposes, multiplied
by five. Set
\begin{equation}\label{eq:precision64}
 b=L+\lceil\log_2N\rceil+2.
\end{equation}
When $b<N$, initialize $n=(2^b,0,0)\in\Z^3$. On command $a$, compute
\begin{equation}\label{eq:integerupdate64}
 z=M_an,\qquad
 n'_i=\operatorname{sgn}(z_i)\lfloor |z_i|/5\rfloor,
 \qquad n\leftarrow n'.
\end{equation}
The current numerical vector is $v=n/2^b$. The update is exactly
$T_b(A_av)$; in particular it does not use an approximate matrix.
The common denominator is kept fixed.

When $b\ge N$, instead initialize $n=e_1$ and $d=1$ and use the exact
updates
\begin{equation}\label{eq:exactupdate64}
                       n\leftarrow M_an,\qquad d\leftarrow5d.
\end{equation}
After $t$ commands, $n/d=A_{a_t}\cdots A_{a_1}e_1$ and $d=5^t$.
No history of the commands is retained in either mode.

\begin{proof}[Upper bound in Theorem~\ref{thm:streamspace64}]
In the fixed-grid mode, Lemma~\ref{lem:directedround64} gives legality
at every step and
\[
 \|D(v_N)-D(x_N)\|_{\rm F}
   =\frac{\norm{v_N-x_N}}{\sqrt2}
   \le\sqrt{3/2}\,N2^{-b}<2^{-L}.
\]
The exact mode is legal and has zero error. In both modes the output
is a rational density matrix. For example, if $v=n/d$, one may output
\begin{equation}\label{eq:output64}
 D(v)=\frac1{2d}
 \begin{pmatrix}d+n_3&n_1-\mathrm i n_2\\
                  n_1+\mathrm i n_2&d-n_3\end{pmatrix}.
\end{equation}
The guarantee $\norm n\le d$ implies positive semidefiniteness.
There is no square root, eigenvalue computation, projection oracle, or
physical quantum preparation in the algorithm.

In the fixed-grid mode each coordinate needs $b+O(1)$ bits, including
its sign. Every entry of $M_a$ has absolute value at most five and each
row has absolute sum at most seven. Thus the temporary integer $z_i$
and the new coordinate also require only $b+O(1)$ bits. A constant
number of buffers suffices. Addition, multiplication by these fixed
integers, and division of a nonnegative integer by five are executable
by binary digit scans with $O(b)$ bit operations and $O(b)$ storage.
Truncation of a negative coordinate uses its sign and absolute value;
rounding toward negative infinity would not give this algorithm.

In the exact mode $\norm n=d=5^t$, so all numerators, the denominator,
and intermediate results have $O(t+1)$ bits. Their updates use
$O(t+1)$ bit operations and storage. Taking $b<N$ for the first mode
and $b\ge N$ for the second yields \eqref{eq:spaceupper64} and the
asserted command time. Formula~\eqref{eq:output64} can be written as
signed binary integers over one common denominator in $O(b_*)$ time
and space. Conversion to a decimal presentation is not needed for the
bit-operation bound.

For completeness, parameter handling does not require storing a very
large requested precision. Read $N$ and store it in binary. While
reading $L$, maintain its value capped at $N$ by the recurrence
$\ell\leftarrow\min\{N,2\ell+\delta\}$ for each binary digit
$\delta$. If the cap is reached, the exact mode is already determined.
Otherwise $L<N$ is stored exactly, and \eqref{eq:precision64} is
computed with $O(\log N)$ auxiliary bits. This also works with any
fixed input radix. Reading the entire header takes its actual input
length; a straightforward implementation uses $O(\log N)$ work bits
and $O(\log N)$ bit operations per header digit. A counter validating
that exactly $N$ commands were read costs $O(\log N)$ additional bits,
which are within the asserted bound. No command-time external clock
is needed.
\end{proof}

\begin{corollary}[The zero-error and arbitrary-tolerance regimes]\label{cor:allprecision64}
Exact one-pass numerical realization of this fixed rational experiment
has bit-space $\Theta(N)$. For every real $0<\varepsilon\le1/4$, a
sufficient binary accuracy parameter $L=\lceil\log_2(1/\varepsilon)\rceil$
gives optimal space order
\[
              \Theta\bigl(\min\{N,\log(N+1)+\log(1/\varepsilon)\}\bigr).
\]
Thus the order is logarithmic at fixed positive error and linear once
$\log(1/\varepsilon)$ is of order $N$ or larger.
\end{corollary}
\begin{proof}
Use the exact profile $5^N$ for the first lower bound and
\eqref{eq:exactupdate64} for its upper bound. For positive tolerance,
the lower proof applies directly with that tolerance in
\eqref{eq:bitdichotomy64}. The chosen $L$ differs from
$\log_2(1/\varepsilon)$ by less than one and supplies the upper bound.
\end{proof}

\subsection{Rational orthogonal actions}
The computational consequence is not restricted to the numerical LPS
example. Legal unit-ball outputs allow the same argument for any fixed
rational orthogonal alphabet satisfying the two geometric hypotheses
used by the entropy theorem.

\begin{corollary}[Uniform space for rational orthogonal actions]\label{cor:rationalspace64}
Let $\mathcal A\subset O(\Q^m)$ be a fixed finite alphabet, let $G$
be its compact generated group, and let $x_0\in\Q^m$ be a unit vector.
Assume the full action gap \eqref{eq:fullgap62} and the all-direction
cap bound \eqref{eq:caps62} with a positive exponent $p$.
Require numerical decoder values to belong to the closed unit ball
of $\R^m$. In the one-pass model, with Euclidean mean error at most
$2^{-L}$, the optimal uniform bit-space order is
\[
             \Theta\bigl(\min\{N,L+\log_2(N+1)\}\bigr)
\]
for $L\ge2$ and $N$ above a constant depending only on the fixed
input. There is no identity or return hypothesis. The upper bound is
deterministic with pointwise error. Exact computation has order $N$.
\end{corollary}
\begin{proof}
The lower proof of Theorem~\ref{thm:driftoccupation62} uses the legal
output set only through $\norm{D_i}\le1$ in the terminal centroid
estimate. It therefore applies unchanged to unit-ball decoders, with
$\rho=1$ and $\zeta=1-\varepsilon$. Its finite budget gives
\[
 N\le C_0+C_1\log k+C_2\varepsilon k^{2/p}
                   \qquad(0\le\varepsilon\le1/4).
\]
The same dichotomy as in \eqref{eq:bitdichotomy64} gives
$\log k\ge c\min\{N,\log(N/\varepsilon)\}-C'$.
For sufficiently large $N$ this lower bound absorbs its additive
constant uniformly, since $\log(N/\varepsilon)\ge\log(4N)$.
At zero error the finite budget directly gives $\log k\ge cN-C'$.
The complete-configuration reduction then applies.

For the upper bound write $A_a=M_a/D$ and $x_0=n_0/d_0$ with fixed
integer data and positive common denominators. At grid precision $b$
initialize $v_0=T_b(x_0)$ and perform the integer update
$n'_i=\operatorname{sgn}((M_an)_i)\lfloor |(M_an)_i|/D\rfloor$.
The initial error contributes one additional rounding term, so
Lemma~\ref{lem:directedround64} gives
$\norm{v_N-x_N}\le(N+1)\sqrt m\,2^{-b}$ and legal outputs at
all prefixes. Choose
$b=L+\lceil\log_2(N+1)\rceil+\lceil\log_2\sqrt m\rceil+2$.
Alternatively retain the exact vector with denominator $d_0D^t$.
Since the dimension and integer coefficients are fixed, these modes
use respectively $O(b)$ and $O(N+1)$ bits, including scratch and
parameter storage. Choose the smaller mode. The capped parameter
scan and fixed-integer digit operations from the preceding proof
apply without change. No orbit-covering algorithm or construction of
a real transition table is used.
\end{proof}
This corollary requires unit-ball decoders. It does not assert that
coordinatewise truncation preserves an arbitrary orbit hull or the
positive-semidefinite cone in higher matrix dimension. The Bloch ball
is exactly the legal qubit output set, which is why the preceding
explicit result gives legal density matrices.

\subsection{Relation to finite automata and positive realization}
The passage from a configuration to a stochastic label in
Lemma~\ref{lem:boundary64} is elementary, as is directed finite-precision
arithmetic. The quantitative content is that the entropy obstruction
continues to force optimal uniform bit-space even when arbitrary
randomization and unlimited internal running time are allowed.
This is not a reinterpretation of free real-table lookup as an
implemented algorithm. The upper program is independently specified.
The logarithm of width records fewer asymptotic details than width
itself: the present matching space order does not remove the
subexponential multiplicative loss in Theorem~\ref{thm:crossover63}.

The invariant-cone characterization of positive realization, stated by
Benvenuti--Farina \cite[Theorem~2]{BF64}, concerns exact stationary
nonnegative realizations of a rational transfer function. Our lower
bound instead concerns all-word, horizon-dependent approximation by
common stochastic rows. The affine Bloch recursion has constant real
linear dimension, but its negative matrix entries do not make it a
positive finite-label realization. A finite-bit implementation of that
recursion must pay for its accuracy.

Ambainis--Watrous \cite[Section~2 and Theorem~1]{AW64} allow a two-way
input head and a fixed-dimensional quantum register in a language
recognition problem. Panduranga Rao--Vinay
\cite[Theorem~1]{RV64} construct complex-weighted two-way automata
for one-sided language recognition. Their transitions are not
nonnegative stochastic rows with a legal density-matrix decoder.
Neither comparison identifies algebraic amplitudes with finitely many
classical configurations. The model here has a one-pass classical
input, no physical quantum register, and a numerical output required
for every command word. We do not claim a language-recognition
separation or an extension to repeated quantum measurements.
